\section{Terms inside a factor: the interaction term}\label{f:ch:interaction}
\sectionbrief
{an evolution factor $u_t=Au$ with $A=A_1+\dots+A_N$, the \emph{terms} (advection, diffusion, dispersion, \dots), each a polynomial $a_n(\partial_x)$; flows $S_n(t)=e^{tA_n}$; a bounded rectangle $\Omega$ and a base time $t_0$. The terms commute, except in Section~\ref{f:sec:noncomm}.}
{Proposition~\ref{f:prop:terms} shows that splitting a solution by terms is never exact. Here we compute exactly what is left over, the interaction term $\Eint$, how much it costs, and whether it can be computed stably.}
{third link of the chain, answering Q3. The answer is positive in a function class of exponential type (an exact product formula, cost $O(N)$ per value) and negative in Sobolev spaces (no continuity, no stable regularisation).}

\subsection{From the sum of operators to the interaction term}

The solution of $u_t=Au$ is given by the flow (semigroup) of $A$:
\[
u(t)=e^{tA}u_0=S(t)\,u_0 .
\]
If term $n$ acted alone, it would move the datum to $S_n(t)u_0=e^{tA_n}u_0$, a change of $S_n(t)u_0-u_0$.

The naive guess is that the terms act independently, so that the total change is the sum of the individual changes. It is not: the flow of a sum is not the sum of the flows. For advection ($A_1=-c\partial_x$, a translation) and diffusion ($A_2=D\partial_x^2$, a smoothing), the solution is translated \emph{and} smoothed. This joint action is not ``only translated'' plus ``only smoothed''.

\begin{definition}[{Interaction term}]\label{def:eint}
With $S_n(t)=e^{(t-t_0)A_n}$ the flow of term $n$ alone,
\begin{equation}\label{f:eq:eint-def}
\Eint^{\mathrm{add}}(t)=u(t)-u_0-\sum_{n=1}^N\big(S_n(t)u_0-u_0\big).
\end{equation}
The multiplicative variant compares the true flow with the product of the single flows:
\[
\Eint^{\mathrm{mult}}(t)=e^{(t-t_0)A}-\prod_{n=1}^Ne^{(t-t_0)A_n}.
\]
We write $\Eint$ for $\Eint^{\mathrm{add}}$ unless the distinction matters.
\end{definition}

\subsection{The function class and the exact formula}

Before computing, one must decide on which functions the flows of the single terms are defined. This is not a formality. For the heat equation there are nonzero solutions with zero initial datum (Tychonoff), so ``the diffusion flow of $u_0$'' is not unique without a class.

\begin{definition}[{The class $\Bclass$ of exponential type \led{P1-D1}}]\label{def:B}
For $\rho>0$, $B_\rho(J)$ is the space of smooth $f$ on the interval $J$ with
\[
\|f\|_{B_\rho}^2=\sum_{j\ge0}\rho^{-2j}\,\|f^{(j)}\|_{L^2(J)}^2<\infty .
\]
$\Bclass_\rho(\Omega)$ is the analogous space in $(t,x)$, and $\Bclass=\bigcup_\rho\Bclass_\rho$. Every exponential-polynomial belongs to $\Bclass$, and on $B_\rho$ the derivative is bounded: $\|\partial_x\|\le\rho$.
\end{definition}

On $\Bclass$ every polynomial in $\partial_x$ is a bounded operator, its exponential is an absolutely convergent series, and two such exponentials commute \led{P1-L1, P1-L2} \proved. So the algebra of exponentials of numbers applies to operators.

\begin{theorem}[{Closed form of the interaction \led{P1-T2} \proved}]\label{f:thm:eint}
Let $u\in\Bclass_\rho(\Omega)$ solve $(\partial_t-a(\partial_x))u=0$ with $a=\sum_na_n$, and put
\[
\Delta_n=e^{(t-t_0)a_n(\partial_x)}-\Id .
\]
Then the single flows are unique in $\Bclass$, $\Eint^{\mathrm{mult}}\equiv0$, and
\begin{equation}\label{f:eq:eint}
\Eint^{\mathrm{add}}(t)=\Big[\prod_{n=1}^N(\Id+\Delta_n)-\Id-\sum_{n=1}^N\Delta_n\Big]u_0
=\sum_{|S|\ge2}\ \prod_{n\in S}\Delta_n\,u_0 .
\end{equation}
For a datum $u_0=\sum_kc_ke^{\zeta_kx}$, each mode is multiplied by the number
\[
e_{\ge2}(z)=\prod_n(1+z_n)-1-\sum_nz_n,\qquad z_n=e^{(t-t_0)a_n(\zeta_k)}-1 .
\]
\dimref{P1-T2}
\end{theorem}

\begin{intuition}
Expand the product $\prod_n(\Id+\Delta_n)$. It gives one monomial per subset $S$ of terms. The empty subset gives $\Id$, the singletons give the individual changes $\Delta_n$, and everything else is interaction. So the interaction is the sum of products of \emph{at least two} increments: it appears only when at least two terms act together.

For two terms the formula is simply
\[
\Eint=\Delta_1\Delta_2\,u_0=(S_1-\Id)(S_2-\Id)\,u_0 .
\]
In each Fourier mode, the interaction is the product of the two single-term responses.
\end{intuition}

\begin{example}[{Advection--diffusion with a Gaussian datum}]\label{ex:advdiff}
Take $u_t=-cu_x+Du_{xx}$ and $u_0=g(x;\sigma^2)$, where
\[
g(x;v)=\frac{e^{-x^2/2v}}{\sqrt{2\pi v}} .
\]
The two flows are a translation and a Gaussian smoothing:
\[
S_1u_0=g(x-ct;\sigma^2),\qquad S_2u_0=g(x;s^2),\qquad u=g(x-ct;s^2),
\]
with $s^2=\sigma^2+2Dt$. The interaction term is
\[
\Eint=g(x-ct;s^2)-g(x-ct;\sigma^2)-g(x;s^2)+g(x;\sigma^2).
\]
Figure~\ref{f:fig:eint}(a) shows that it is comparable in size to $u$ itself: it is not a small correction. In Fourier variables
\[
\frac{|\widehat{\Eint}(\kappa)|}{|\hat u_0(\kappa)|}=\big|e^{-ic\kappa t}-1\big|\cdot\big|e^{-D\kappa^2t}-1\big|\le2 .
\]
It vanishes at $\kappa=0$ and whenever $c\kappa t\in2\pi\Z$, where the translation returns to the identity.
\end{example}

\begin{figure}[!htb]
\centering
\begin{tikzpicture}
\begin{axis}[width=7.4cm,height=5.6cm,xlabel={$x$},xmin=-6,xmax=8,ymin=-0.2,ymax=0.45,
 title={\small (a) $t=1$, $c=2$, $D=\tfrac12$},legend style={font=\scriptsize,draw=none,fill=white,fill opacity=0.8,text opacity=1,at={(0.02,0.98)},anchor=north west},
 domain=-6:8,samples=140,no markers]
\addplot[cG,thick,dashed] {gauss(x,1)};
\addplot[cA,thick] {gauss(x-2,2)};
\addplot[cB,very thick] {gauss(x-2,2)-gauss(x-2,1)-gauss(x,2)+gauss(x,1)};
\legend{$u_0$,$u(t)$,$\Eint$}
\end{axis}
\end{tikzpicture}\hfill
\begin{tikzpicture}
\begin{axis}[width=7.0cm,height=5.6cm,xlabel={$\kappa$},xmin=0,xmax=4,ymin=0,ymax=2.1,
 title={\small (b) $|\widehat{\Eint}|/|\hat u_0|$},legend style={font=\scriptsize,draw=none,at={(0.5,-0.28)},anchor=north,legend columns=3},
 domain=0:4,samples=160,no markers]
\addplot[cA,thick] {sqrt(((cos(deg(2*0.5*x))-1)^2+sin(deg(2*0.5*x))^2)*(exp(-0.5*0.5*x^2)-1)^2)};
\addplot[cC,thick] {sqrt(((cos(deg(2*1.0*x))-1)^2+sin(deg(2*1.0*x))^2)*(exp(-0.5*1.0*x^2)-1)^2)};
\addplot[cB,thick] {sqrt(((cos(deg(2*2.0*x))-1)^2+sin(deg(2*2.0*x))^2)*(exp(-0.5*2.0*x^2)-1)^2)};
\legend{$t=0.5$,$t=1$,$t=2$}
\end{axis}
\end{tikzpicture}
\caption{Advection--diffusion with Gaussian datum ($\sigma^2=1$). (a) The datum, the solution and the interaction term. (b) The Fourier ratio of the interaction to the datum, for three times.}\label{f:fig:eint}
\end{figure}

\begin{example}[{A stationary state is pure interaction}]
For advection--diffusion, a stationary solution satisfies $-cu_x+Du_{xx}=0$, for example $u=e^{cx/D}$. The full flow leaves it unchanged, so $u(t)-u_0=0$. The single flows do change it, and the interaction term exactly cancels their changes. In this mode
\[
\Eint=-4\sinh^2\!\Big(\frac{\theta}{2}\Big)\,u,\qquad \theta=\frac{c^2(t-t_0)}{D}.
\]
Equilibrium is the case where all the dynamics is interaction.
\end{example}

\subsection{Cost: $O(N)$ versus $2^N$}

Expanding~\eqref{f:eq:eint} produces one monomial per subset with at least two elements: $2^N-N-1$ cross terms (Figure~\ref{f:fig:lattice}). Writing the \emph{complete table} of cross terms is therefore exponential. Computing one \emph{value} is cheap, because the product is never expanded.

\begin{figure}[!htb]
\centering
\begin{tikzpicture}[x=1.8cm,y=1.1cm,every node/.style={font=\small}]
\node (e) at (0,0) {$\emptyset$};
\node (1) at (-1,1) {$\{1\}$}; \node (2) at (0,1) {$\{2\}$}; \node (3) at (1,1) {$\{3\}$};
\node[text=cB] (12) at (-1,2) {$\{1,2\}$}; \node[text=cB] (13) at (0,2) {$\{1,3\}$}; \node[text=cB] (23) at (1,2) {$\{2,3\}$};
\node[text=cB] (123) at (0,3) {$\{1,2,3\}$};
\foreach \a/\b in {e/1,e/2,e/3,1/12,1/13,2/12,2/23,3/13,3/23,12/123,13/123,23/123} \draw[cG] (\a)--(\b);
\node[anchor=west,align=left,font=\scriptsize] at (1.8,2.2) {\textcolor{cB}{red}: cross terms $\prod_{n\in S}\Delta_n$, $|S|\ge2$\\ there are $2^N-N-1$ of them};
\end{tikzpicture}
\caption{The subset lattice for $N=3$ terms. Each node with at least two elements is one cross term of the interaction.}\label{f:fig:lattice}
\end{figure}

\begin{theorem}[{Cost of the interaction \led{P1-T3} \proved}]\label{thm:cost}
Let the datum be a sum of $M$ exponentials with polynomial prefactors of degree at most $L$, and let $d_n=\deg a_n$.
\begin{enumerate}[label=(\alph*)]
\item One value of $\Eint$ costs
\[
O\Big(M\big[N(L+1)^2+(L+1)\textstyle\sum_n(d_n+1)\big]\Big);
\]
for simple exponentials ($L=0$) this is $O\big(M(N+\sum_nd_n)\big)$: linear in the number of terms.
\item All aggregates by order, $E^{(k)}=\sum_{|S|=k}\prod_{n\in S}\Delta_nu_0$ for $2\le k\le N$, cost $O(MN^2(L+1)^2)$ plus the same term in $d_n$.
\item The full table of cross terms has $M(2^N-N-1)$ entries and costs $\Theta\big(M\,2^N(L+1)^2\big)$.
\end{enumerate}
\dimref{P1-T3}
\end{theorem}

\begin{example}[{Counting operations for one mode}]
For one simple mode $e^{\zeta x}$ and $N$ terms, compute the $N$ numbers $z_n=e^{(t-t_0)a_n(\zeta)}-1$, then the product $\prod_n(1+z_n)$ with $N-1$ multiplications, then subtract $1+\sum_nz_n$. That is about $3N$ operations. Listing the $2^N-N-1$ products $\prod_{n\in S}z_n$ one by one costs $2^N$ times more for large $N$: already for $N=20$, about a million.
\end{example}

\subsection{The flows are not determined without a class}

\begin{proposition}[{Non-uniqueness of single flows \led{P1-P1(a)--(d)} \proved}]\label{prop:nonuniq}
On a bounded rectangle, without a function class:
\begin{enumerate}[label=(\alph*)]
\item for the transport term $a_n=\xi$, the flow of a datum is any smooth extension of it: infinitely many choices;
\item for the diffusion term $a_n=\xi^2$, there are nonzero heat solutions with zero datum (Tychonoff) \cf;
\item so $\Eint$ is not a function of $u$ alone; inside $\Bclass$ it is unique.
\end{enumerate}
\dimref{P1-P1}
\end{proposition}

\subsection{Stability: the interaction cannot be computed stably}

\begin{theorem}[{No continuous extension to Sobolev spaces \led{P1-P1(e)} \falsest}]\label{f:thm:nocont}
Let $Q=\tau+c\xi-D\xi^2$ with $c,D>0$ (advection--diffusion), terms $a_1=-c\xi$ and $a_2=D\xi^2$, and any base time $t_0$. For every $m$,
\[
\sup\Big\{\frac{\|\Eint\|_{L^2(\Omega)}}{\|u\|_{H^m(\Omega)}}:\ u\in\EP,\ Qu=0,\ u\ne0\Big\}=\infty .
\]
So $u\mapsto\Eint$ does not extend continuously from exponential-type data to $H^m$, although exponential-type data are dense. The same holds for every evolution with at least two non-proportional terms; for proportional terms the interaction is continuous \led{OPA-P3} \proved.

\dimref{P1-P1e}
\end{theorem}

\begin{intuition}
On a bounded interval, a real exponential $e^{\zeta x}$ with large $\zeta$ is a legitimate datum. The diffusion flow alone amplifies it by $e^{D\zeta^2t}$. The true solution is only changed by the factor $e^{(D\zeta^2-c\zeta)t}$ times the datum norm. The interaction term must cancel the gap, and the gap grows faster than any Sobolev norm of the datum.
\end{intuition}

This instability has three consequences for any algorithm that wants to be stable.

\begin{theorem}[{No regularisation removes the interaction \led{P2-T1, P2-T2, P2-T4} \proved}]\label{f:thm:noreg}
\begin{enumerate}[label=(\alph*)]
\item \textbf{Any penalty leaves the same residual.} For any regulariser (an $\ell^1$ sparsity penalty, $\ell^2$, total variation, \dots) the residual of the best fit converges to the distance of $u$ from the closed span of the per-term pieces.
\item \textbf{Inside an evolution factor that residual is large.} If no single term exhausts the symbol, the part of $u$ outside the span of the per-term kernels satisfies
\[
\|E^\perp(u)\|_{H^m}\ \ge\ \frac{\|g(\partial_x)u\|_{H^{m-k}}}{C},\qquad g=\prod_n(a-a_n).
\]
Only a finite-dimensional space of exponential-polynomials has a per-term representation. For advection--diffusion it is
\[
\spn\{1,\ x-ct,\ (x-ct)^2+2Dt\}.
\]
\item \textbf{No uniformly Lipschitz regularisation exists.} Any family of Lipschitz maps converging pointwise to an unbounded map has Lipschitz constants tending to infinity. Every exact selection of components has Lipschitz constant at least $\gamma_N^{-1/2}$.
\end{enumerate}
\dimref{P2-T2}
\end{theorem}

\begin{caution}[title=Added in the complete edition (Phase 10)]
The instability above is proved for advection--diffusion. The general statement ``for every operator with at least two non-proportional terms'' (P1-Q6) is false as written: with constant terms, for example reaction--diffusion, the interaction is continuous from $H^1$. The correct dichotomy is by the number of terms of degree $\ge1$ (Theorem~\ref{thm:s10at1}, Section~\ref{ch:phase10}).
\end{caution}

\begin{aside}{The one continuous split found}
A \emph{balance} decomposition of $u$ writes each term's contribution as the time integral of that term applied to the true solution:
\[
u^{(i)}(t)=\int_{t_0}^{t}a_i(\partial_x)\,u(s)\dd s .
\]
It is exact ($u-u_0=\sum_iu^{(i)}$) and continuous in $H^m$ \led{P3-T6} \pw. But it needs an external principle to decide that this is ``the'' contribution of term $i$: it uses the full solution, not the flow of the term alone.
\end{aside}

\subsection{Time-dependent perturbation theory}

The product formula needs commuting terms. With non-commuting terms, or with boundary conditions, there is no closed formula, but $\Eint$ is still defined by~\eqref{f:eq:eint-def}. Perturbation theory expands the flow in powers of the interaction. Split $A=A_1+B$, where $A_1$ is the principal term and $B$ the rest. The Dyson--Phillips series is
\begin{equation}\label{f:eq:dyson}
\begin{aligned}
S(t)&=\sum_{k\ge0}S^{(k)}(t),\\
S^{(0)}(t)&=e^{tA_1},\\
S^{(k+1)}(t)&=\int_0^te^{(t-s)A_1}\,B\,S^{(k)}(s)\dd s .
\end{aligned}
\end{equation}
$S^{(k)}$ collects the histories in which $B$ acts $k$ times, at intermediate times, with the flow of $A_1$ in between. This is the interaction picture of quantum mechanics~\cite{dy}.

\begin{theorem}[{Dyson--Phillips converges for every time \led{P3-T1} \proved}]\label{thm:dyson}
Let $A_1=a(\partial_x)$ with symbol of degree $q$, and $B=\sum_{k\le r}b_k(x)\partial_x^k$ with bounded coefficients and $r<q$. The symbol condition
\[
\mathrm{Re}\,a(i\kappa)\le\omega-\delta|\kappa|^q\qquad(\delta>0)
\]
holds if and only if $q$ is even and the leading coefficient has negative real part (a dissipative operator of even order). Under it, the Dyson--Phillips series converges in operator norm, uniformly on every bounded time interval, and its sum is the semigroup generated by $A_1+B$. The same holds with Dirichlet conditions on an interval.

\dimref{P3-T1}
\end{theorem}

So $\Eint$ is well defined and can be approximated by truncating the series even without a closed formula.

\begin{proposition}[{Interaction created by a boundary \led{P1-P6, P3-T3, OPA-T2, OPF5A-T5} \proved}]\label{prop:boundary}
On an interval with Dirichlet conditions, advection and diffusion no longer commute: the boundary creates an $\Eint^{\mathrm{mult}}$ already at first order \led{P1-P6}. For small times its size is governed by how the datum touches the boundary. If $k$ is the order of contact (the first nonvanishing derivative of $u_0$ at the endpoints), then
\[
\|\Eint^{\mathrm{mult}}(t)\|=\Theta\big(t^{k/2+3/4}\big)\qquad(t\to0).
\]
For compatible data the leading constant is explicit and proportional to
\[
|c|\,\frac{\tilde N_k}{k!}\,D^{k/2-1/4}\sqrt{u_0^{(k)}(0)^2+u_0^{(k)}(L)^2}.
\]
\dimref{OPF5A-T5}
\end{proposition}

The smoother the datum is against the boundary, the faster the boundary interaction becomes negligible as $t\to0$.

\subsection{Non-commuting terms}\label{f:sec:noncomm}

If two terms do not commute, the flow of their sum is not the product of the flows. By the Baker--Campbell--Hausdorff formula,
\[
e^{tA_1}e^{tA_2}=\exp\Big(t(A_1+A_2)+\frac{t^2}{2}[A_1,A_2]+O(t^3)\Big).
\]
The commutators generate a Lie algebra, and its type decides what can be done.

\begin{theorem}[{Lie-algebra classification \led{P1-T4, P1-P5, P1-T5} \proved}]\label{f:thm:lie}
\begin{enumerate}[label=(\alph*)]
\item \textbf{Finite and solvable.} An exact formula exists, with $d=\dim\mathfrak g$ entire scalar functions (Wei--Norman). Example: $\partial_x^2+x$, and the Ornstein--Uhlenbeck operator.
\item \textbf{Finite, not solvable} (for example $\mathfrak{sl}_2$). A closed form exists locally in time; it can fail globally.
\item \textbf{Infinite adjoint orbit.} If $X$ and all the iterated commutators $\mathrm{ad}_X^kY$ are linearly independent (for example $X=x^a\partial_x$, $Y=\partial_x$ with $a\ge3$), then no finite product of flows of the two terms, with times analytic in $t$ and in the coupling $\varepsilon$ and with the flows of $Y$ switched off at $\varepsilon=0$, reproduces $e^{t(X+\varepsilon Y)}$. Only series expansions exist.
\end{enumerate}
\dimref{P1-T5}
\end{theorem}

\begin{example}[{Which potentials allow Wei--Norman}]
For $i\partial_t=-\partial_x^2+V(x)$, the operators $\partial_x^2$ and $x^2$ generate the finite Lie algebra $\mathfrak{sl}_2$: the harmonic oscillator has a closed-form evolution. With $V=x^3$, repeated commutators produce polynomials of ever higher degree and the algebra is infinite-dimensional. But the hypothesis of part (c) fails there, so whether a finite product of flows is impossible for $\partial_x^2$ and $x^k$ with $k\ge3$ is still \openq{} \led{P1-Q3}. Part~(c) applies to the transport with polynomial speed $u_t=(x^a+\varepsilon)u_x$, $a\ge3$, which nevertheless has an explicit solution by characteristics: the theorem excludes only a factorisation into flows of the terms.
\end{example}

\begin{caution}[title={Other ways commutativity is lost}]
\begin{itemize}
\item \textbf{Repeated factors.} For $(\tau-a)^e$ the ``multiplicity term'' is a nilpotent shift. It commutes with the other terms, its flow is polynomial in time (secular growth), and $\Eint^{\mathrm{mult}}=0$ \led{OPF-T2} \pw.
\item \textbf{Second-order evolutions} (wave, Klein--Gordon). The companion blocks do not commute, so the first-order calculus fails; a Gauss factorisation works mode by mode \led{II-T7, P1-P4} \proved.
\end{itemize}
\end{caution}

\begin{keypoint}
For commuting terms the interaction is explicit and cheap to evaluate, but it can be neither removed nor computed stably in Sobolev spaces. That instability is not a defect of the formula: it is a property of the question. For non-commuting terms even the explicit form is lost, except when the terms generate a finite solvable Lie algebra.
\end{keypoint}

